% ===================== 第 1 章 =====================
\section{研究背景与设计立场}

\subsection{深度研究管线的三个结构性缺陷}

把一个争议问题交给单管线深度研究智能体，它通常能交出一篇行文流畅、引用齐全的综述。流畅本身没有问题，问题在于争议问题要求的不是流畅，而是把证据之间的裂缝暴露出来。标准管线在这件事上有三个先天缺陷，且都无法靠调整提示词消除。

\textbf{第一，视角由管线结构决定，缺了就是缺了。}一条没有测量学家角色的管线，不会主动质疑「社交媒体使用强度」这种自报告量表的构念效度；一条没有因果推断角色的管线，会把横截面相关顺畅地写成因果。视角缺失不会报错，它只会沉默地体现在结论里。

\textbf{第二，重复不等于证实。}如果多个智能体检索到同一篇被撤稿的论文、同一条张冠李戴的引用、同一份只发了摘要的工作，系统可以让它们彼此「印证」，产出一个看起来被多源验证、实则只有一个错误来源的共识。论文数量与独立证据数量是两件事，普通管线不区分。

\textbf{第三，总结步骤天然倾向收敛。}总结器的优化目标是一段通顺文字，于是少数派反对什么、需要什么证据才会改判，这些让结论「不光滑」的信息最先被牺牲。争议研究的价值恰恰在这些不光滑处。

\subsection{领域背景：为什么先做计算社会科学}

这些风险在计算社会科学领域尤其具体。心理学的可复现性评估显示，大量已发表效应在重复实验中显著缩水或消失\footnote{Open Science Collaboration, \emph{Estimating the reproducibility of psychological science}, Science, 2015.}；发表偏差系统性地夸大表观效应量\footnote{D. Fanelli, ``Positive'' results increase down the hierarchy of the sciences, PLoS ONE, 2010.}；研究者自由度（选择变量、切分样本、挑选模型规格）会成倍放大假阳性\footnote{A. Gelman \& E. Loken, \emph{The garden of forking paths}, 2013；B. Nosek et al., \emph{The preregistration revolution}, PNAS, 2018.}。

在首版聚焦的数字行为与心理健康方向，问题更加密集：横截面研究被例行地叙述成因果，自报告测量占据主流，少数几个公开数据集支撑着大量彼此并不独立的文献\footnote{A. Orben \& A. Przybylski, \emph{The association between adolescent well-being and digital technology use}, Nature Human Behaviour, 2019；P. Valkenburg et al., \emph{Social media use and its impact on adolescent mental health}, Current Opinion in Psychology, 2022.}。这是一个「盲点发现」边际价值最高的领域，因此 Poliscope 首版明确不做通用全学科，而把这一个领域做深（范围约束见第 12 章）。

\subsection{设计立场：先定组织结构，再放角色}

常见的多智能体研究系统走的是「加人」路线：让若干 Agent 轮流发言，再让一个 Agent 写总结。Poliscope 反其道而行：先回答「一场可靠的争议审查在组织结构上必须满足什么」，再把七个角色放进这套结构。设计立场浓缩为下面三条，它们贯穿后续全部章节，也是评审本系统任何设计决策的标尺。

\begin{principlebox}
\boxtag{\color{accentdk}三条设计原则}
\textbf{1. 异议靠构造保留，不靠主持人自觉。}少数意见、被反驳的观点、未解决的质询都有正式的数据结构承载（\code{DebateCapsule}、\code{DissentCertificate}），总结步骤无权删除它们。\\[3pt]
\textbf{2. 证据靠构造审计，不靠引用格式正确。}每条正式发现都要经过来源、引用蕴含、方法质量三层审计；证据分 A–D 级，等级不够就进不了高置信结论。\\[3pt]
\textbf{3. 过程靠构造隔离，不允许自动变成证据。}科学家在结构上没有证据图写权限；思维链、工具记录、失败路线属于过程图，与证据图物理分离。
\end{principlebox}

\subsection{贡献清单}

\begin{enumerate}[leftmargin=1.6em, itemsep=4pt]
\item \textbf{方法层}：七人科学议会协议——八阶段、八类结构化动作、禁止多数投票、条件化共识与具名异议（第 2、4 章）。
\item \textbf{治理层}：三层记忆与双图证据治理——追加写事件账本、唯一投影器、数据库权限强制写隔离、辩证折叠保留反例（第 5、6、7 章）。
\item \textbf{工程层}：七人全程参与下的成本与可靠性方案——分层模型调用、共享检索缓存、四维预算、单席失败降级、断点续跑、有界实时流（第 8、9、11 章）。
\item \textbf{评测层}：ForesightBlindspot 时间切片基准、五条基线与六项消融，连同真实模型实验的负面结果一起公开（第 10 章）。
\item \textbf{产品层}：以证据地图为中心的 Web 工作台，支持账户体系、个人知识库、研究者技能安装、公开分享与多格式导出（第 11 章）。
\end{enumerate}

\subsection{文档约定}

文内实现陈述均以仓库当前代码为准，引用到模块与符号时使用 \code{packages/}\,/\,\code{apps/} 下的文件路径或类名，便于直接对照。凡是设计规格中保留、首版尚未完整实现的部分，统一在第 12 章集中说明，正文遇到时也会就地标注，不把占位实现写成已完成能力。
